\documentclass[11pt,letterpaper]{article}
\usepackage{luxcover}
\usepackage[margin=1in]{geometry}
\usepackage{amsmath}
\usepackage{booktabs}
\usepackage{longtable}
\usepackage{enumitem}
\usepackage{parskip}
\usepackage{titlesec}
\usepackage{hyperref}

\titleformat{\section}{\Large\bfseries\sffamily}{\thesection}{1em}{}
\titleformat{\subsection}{\large\bfseries\sffamily}{\thesubsection}{1em}{}
\titleformat{\subsubsection}{\normalsize\bfseries\sffamily}{\thesubsubsection}{1em}{}

\hypersetup{colorlinks=true,linkcolor=blue,urlcolor=blue}

\setlist{itemsep=2pt,topsep=4pt}

\title{\textbf{The Post-Quantum Migration Trap}\\[6pt]
\large Lux Industries' PQ Patent Family and the Future Infringement Landscape}
\author{Lux Industries Inc.\\[2pt]\small Internal strategic memorandum}
\date{}

\begin{document}
\luxcoverpage
\maketitle

\section*{Abstract}

In late 2025, Lux Industries Inc. filed a coordinated family of seven post-quantum (PQ) patent applications --- one continuation-in-part of US 11,803,853 B2 (the master CIP) and six focused provisionals --- covering PQ-authenticated validation in blockchain transaction-control systems, PQ data-availability attestation for Layer-2 and Layer-3 scaling, PQ-authenticated cross-chain message gating, PQ institutional tokenization settlement, a cryptographically-agile validation framework, PQ threshold multi-party-computation for regulated custody, and PQ-augmented consensus for Layer-1 validator sets. This paper analyzes the strategic position created by the patent family vis-a-vis future infringers --- specifically, blockchain operators who migrate to post-quantum cryptography in response to NIST standardization (FIPS 203/204/205 in 2024 and HQC selection in 2025) and the corresponding regulatory deadlines. The paper argues that the PQ patent family creates a structural trap: every infrastructure operator forced into PQ migration by regulatory or commercial pressure enters Lux's patent zone, regardless of which specific PQ primitive they choose. The patent family is designed to be primitive-agnostic, jurisdiction-agnostic, and architecture-agnostic across all known blockchain transaction-control surfaces, providing Lux Industries with a moat extending the original Nahmii patent through 2046 and beyond.

\section{The post-quantum migration is mandatory}

\subsection{NIST standardization and regulatory pressure}

The post-quantum cryptographic migration is no longer optional. NIST finalized three core PQC standards in 2024: FIPS 203 (Module-Lattice Key-Encapsulation Mechanism, ML-KEM); FIPS 204 (Module-Lattice Digital Signature Algorithm, ML-DSA); FIPS 205 (Stateless Hash-Based Signature, SLH-DSA). NIST selected Hamming Quasi-Cyclic (HQC) as an additional backup KEM in 2025. NIST's stated position is that these algorithms ``can and should be put into use now.''

Comparable mandates have been issued by:

\begin{itemize}
\item European Telecommunications Standards Institute (ETSI) --- PQ migration timelines for telecom and financial infrastructure.
\item European Union Agency for Cybersecurity (ENISA) --- PQ recommendations for critical infrastructure.
\item Bundesamt für Sicherheit in der Informationstechnik (BSI, Germany) --- mandates PQ migration with specific timelines for federal systems and recommends migration for industry.
\item United Kingdom National Cyber Security Centre (NCSC) --- PQ migration guidance for UK critical national infrastructure.
\item Australian Signals Directorate (ASD) --- PQ migration mandates for Australian government systems.
\item Monetary Authority of Singapore (MAS) --- PQ readiness expectations for licensed financial institutions.
\end{itemize}

Beyond regulatory mandate, the commercial logic is forcing: store-now-decrypt-later attacks are already happening against classical-encrypted financial data. Any institutional financial actor with multi-decade data-protection obligations is implementing or planning PQ migration on a 2-5 year horizon.

\subsection{The migration cannot avoid blockchain transaction-control surfaces}

A first-order observation: every blockchain system that authenticates validation results with classical signatures (which is, today, essentially every blockchain system) must migrate. The migration touches:

\begin{itemize}
\item Validator quorum certificates on Layer-1 chains.
\item Rollup operator-and-challenger attestations on Layer-2 chains.
\item Data-availability committee signatures on Alt-DA and DAC chains.
\item Decentralized verifier network (DVN) and DON attestations on cross-chain messaging protocols.
\item Validator signatures on bridges and interoperability hubs.
\item Threshold MPC signatures for institutional custody.
\item Compliance and transfer-agent attestations on regulated tokenization stacks.
\item Light-client signature verification for embedded blockchain consumers.
\end{itemize}

Every one of these surfaces is in Lux's existing US 11,803,853 patent zone (transaction control with operator + verifier + execution gate). The PQ family extends the zone to cover the PQ-authenticated variants.

\section{The PQ patent family}

The PQ family comprises seven applications coordinated as a single prosecution program, filed concurrently to establish a unified priority date for the new matter.

\subsection{Master continuation-in-part: 2026-001}

The umbrella application is a CIP of the parent. It claims an architecture in which validation attestations gating execution of blockchain transactions are authenticated using post-quantum or hybrid classical-and-post-quantum signature primitives, with explicit support for threshold attestation, primitive agility, and precompiled-contract on-chain verification. The CIP inherits the parent's 4 May 2018 priority date for any subject matter present in the parent specification, and claims a new priority date for the PQ-specific matter.

Independent claim 1 of the CIP captures the structural requirement: \emph{validation attestations produced by the second set of validation nodes are post-quantum authenticated; the second validation agent obtains a threshold set of post-quantum authenticated attestations; the first execution agent permits or denies execution in response to verification of the threshold post-quantum attestation set.} Dependent claims cover specific primitives (ML-DSA, SLH-DSA, FN-DSA, hybrid constructions), threshold constructions (bare-threshold, DKG, SNARK-of-signatures), and embodiments per architectural surface.

\subsection{Six focused provisionals: 2026-002 through 2026-007}

Each provisional targets a specific architectural surface and provides claim-level fidelity to that surface's distinctive primitives:

\begin{longtable}{lp{0.7\linewidth}}
\toprule
\textbf{ID} & \textbf{Title and surface} \\
\midrule
2026-002 & \emph{PQ Threshold Data-Availability Attestation.} Targets Alt-DA / DAC / AnyTrust architectures in rollups and L3 chains. \\
2026-003 & \emph{PQ-Authenticated Cross-Chain Message Gating.} Targets bridges, omnichain messaging, oracle networks, validator-bridge committees. \\
2026-004 & \emph{PQ Institutional Tokenization Settlement.} Targets regulated finance: ATS/BD/TA stacks, tokenized money-market funds, regulated stablecoins, CBDCs. \\
2026-005 & \emph{Cryptographically Agile Validation Layer.} Targets the algorithm-agility framework: epoch-based primitive rotation, quantum-threat-status triggers, regulatory-directive policy import. \\
2026-006 & \emph{PQ Threshold MPC for Regulated Custody.} Targets the transfer-agent-as-cryptographic-constituent custody pattern. \\
2026-007 & \emph{PQ Consensus with Parallel Classical and PQ Certificates for L1 Validator Sets.} Targets the L1 consensus architecture supporting parallel BLS + ML-DSA + SLH-DSA certificates, strict-PQ profile, light-client matrix exposure. \\
\bottomrule
\end{longtable}

The combination is exhaustive: every commercially-significant PQ-blockchain surface is captured by at least one application in the family.

\subsection{Why a family rather than a single patent}

Three structural reasons drive the family decomposition rather than a single mega-patent.

\textbf{Continuation flexibility.} Each provisional, once converted to non-provisional, generates its own continuation chain. The family thus produces $7 \times 2$ to $7 \times 5$ continuation applications over the prosecution life. Each continuation tightens claim language to read on specific defendants while remaining valid over the prior art. The continuation chain is the prosecutor's natural-language tracker against the evolving infringement landscape.

\textbf{Invalidity isolation.} An IPR against any one application is an IPR against one application only. The defendant must invalidate every application in the family to escape the moat. The cost of comprehensive invalidation is approximately the number-of-applications multiple of single-IPR cost; with seven applications and 30-40\% single-IPR success rate, the probability of total moat collapse is below 1\%.

\textbf{Jurisdictional optimization.} Different provisional families pass §101 / Art. 52 / equivalent subject-matter eligibility tests differently across jurisdictions. Drafting per surface allows tailoring of subject-matter framing per jurisdiction without disrupting other surfaces. The DA-attestation surface, for example, has cleaner technical-effect framing for EPO purposes than the institutional-settlement surface; drafting them separately enables tailored EP prosecution.

\section{The post-quantum migration trap}

The strategic structure of the family is what we term the \emph{post-quantum migration trap}: any blockchain operator forced into PQ migration by NIST mandate, regulatory directive, or commercial pressure enters Lux Industries' patent zone, with no architectural escape route.

\subsection{The trap geometry}

A blockchain operator faces five possible response paths to the PQ migration imperative:

\begin{enumerate}
\item \textbf{Adopt ML-DSA.} Covered by the CIP, the DA-attestation provisional, the cross-chain provisional, the custody provisional, and the consensus provisional. Trap closed.
\item \textbf{Adopt SLH-DSA.} Covered by the same applications. Trap closed.
\item \textbf{Adopt FN-DSA / Falcon.} Covered. Trap closed.
\item \textbf{Adopt hybrid classical-plus-PQ.} Explicitly claimed in every application in the family. Trap closed.
\item \textbf{Maintain classical-only (refuse migration).} Eventually broken by quantum attack; regulatory and commercial pressure forces eventual migration. Trap delayed but not avoided.
\end{enumerate}

Note that the patent family does not specify a single PQ primitive; it claims the \emph{integration of any post-quantum or hybrid signature primitive into the operator-plus-validator-plus-execution-gate architecture}. The trap geometry is therefore primitive-agnostic. Any PQ migration of the architecture infringes; the choice of primitive determines which dependent claim is read, not whether the patent reads.

\subsection{The crypto-agility provisional closes the last loophole}

A sophisticated infringer might attempt to avoid the trap by deploying a non-NIST primitive: an academic lattice scheme, a custom isogeny-based signature, a code-based scheme, a future NIST round-4 selection. Provisional 2026-005 (cryptographically agile validation layer) claims the algorithm-agility framework itself: a policy controller specifying per-epoch active primitives, regulatory-directive policy import, graceful-degradation under cryptanalysis advances, multi-primitive grace periods. Any infrastructure operator deploying a non-standard or future PQ primitive must implement an agility framework to switch in the event of standardization or compromise; that framework infringes provisional 2026-005 regardless of which specific primitives are deployed.

The agility framework is the meta-trap: even avoiding the primitive-specific claims requires implementing the agility framework, which itself is patented.

\subsection{The consensus provisional closes the L1 escape route}

A second sophisticated escape attempt: build an L1 with native PQ consensus from genesis (avoiding migration entirely). Provisional 2026-007 claims the parallel-classical-plus-PQ-certificate architecture (Quasar-class consensus) and the strict-PQ profile. Either profile choice infringes: the parallel mode is the typical migration architecture; the strict-PQ mode is the only viable design for a genesis-PQ chain that wants to drop classical backwards compatibility.

The L1 escape route is closed by the same family.

\section{Future-infringer enumeration}

We enumerate likely future infringers across categories. The list is illustrative, not exhaustive.

\subsection{Layer-1 chains migrating to PQ consensus}

\begin{longtable}{lp{0.6\linewidth}}
\toprule
\textbf{Likely future infringer} & \textbf{Why} \\
\midrule
Ethereum (post-Merge PoS) & Validator set uses BLS aggregate; PQ migration on Ethereum roadmap for late 2020s / early 2030s. Migration architecture is the parallel-classical-plus-PQ pattern of provisional 2026-007 with high probability. \\
Avalanche (P-Chain BLS) & Same BLS-aggregate pattern; same migration target. \\
Solana (Ed25519) & PoH-plus-BFT with Ed25519 validator signatures; PQ migration in Solana research roadmap. \\
Cosmos / Tendermint chains (Ed25519) & Same pattern. \\
Polkadot / Substrate chains (sr25519) & Same pattern. \\
Sui / Aptos (Ed25519) & Same pattern. \\
Bitcoin (ECDSA secp256k1) & ECDSA in scripts; PQ migration under BIP discussion. Address-level migration less directly in patent zone but multisig and Lightning Network upgrades are in zone. \\
\bottomrule
\end{longtable}

\subsection{Layer-2 rollups adopting PQ DA committees or PQ challenger networks}

Every Alt-DA architecture (Mantle, Plasma chains, OP-Stack-with-EigenDA, OP-Stack-with-Celestia, OP-Stack-with-Avail, OP-Stack-with-S3, AltLayer Beacon Layer, Caldera + alt-DA, Conduit + alt-DA, every Arbitrum Orbit AnyTrust chain). All such systems will migrate the DA-committee signature scheme to PQ. Provisional 2026-002 covers the PQ DA-attestation pattern.

Every optimistic rollup adopting PQ-authenticated fault-proof attestations (which all will, eventually). Provisional 2026-002 (in the correctness/fraud variant) and the master CIP cover.

\subsection{Cross-chain messaging migrating to PQ verifier networks}

\begin{longtable}{lp{0.6\linewidth}}
\toprule
\textbf{Protocol} & \textbf{Migration target} \\
\midrule
Chainlink CCIP & Commit DON BLS signatures and RMN ECDSA signatures both migrate to PQ. \\
LayerZero V2 & DVN signature schemes migrate to PQ across the DVN operator base. \\
Axelar & Validator BLS / Ed25519 signatures migrate to PQ. \\
Wormhole & Guardian ECDSA signatures migrate to PQ. \\
Hyperlane & ISM signature schemes migrate to PQ. \\
Cosmos IBC (Tendermint light clients) & Ed25519 PQ migration. \\
Polygon Agglayer & Multi-prover PQ migration. \\
Across, Stargate, Squid, Synapse, deBridge & Each migrates its specific validator-signature scheme. \\
Standalone-bridge competitors (Connext, Sygma, etc.) & Each migrates. \\
\bottomrule
\end{longtable}

Every one of these is captured by provisional 2026-003.

\subsection{Institutional tokenization migrating to PQ-compliant transfer-agent and custody}

\begin{longtable}{lp{0.6\linewidth}}
\toprule
\textbf{Institution} & \textbf{Migration driver} \\
\midrule
BlackRock (BUIDL, future tokenized portfolios) & SEC and OCC PQ guidance for regulated financial systems. \\
Fidelity Digital Assets & Same. \\
Franklin Templeton (BENJI) & Same. \\
JPMorgan Onyx (Coin Systems, Liink, Token Services) & Same plus internal cryptographic-modernization mandate. \\
Goldman Sachs (DAP) & Same. \\
Citi Token Services & Same. \\
ANZ A\$DC, USDC, USDT, PYUSD issuers & Same. \\
Central bank digital currency projects (ECB digital euro, Fed FedNow extensions, PBOC e-CNY, MAS digital SGD) & Sovereign PQ mandates. \\
DTCC, Euroclear, Clearstream tokenization initiatives & Critical-infrastructure PQ mandates. \\
\bottomrule
\end{longtable}

Every one of these is captured by provisional 2026-004 and provisional 2026-006.

\section{Tactical positioning vs future infringers}

The PQ family creates a structural negotiating position different from the original Nahmii patent. The original patent's value derives from \emph{current} infringement; the PQ family's value derives from \emph{forced future migration}. This changes the negotiation dynamic.

\subsection{The pre-migration license is the high-value moment}

Pre-migration, an infrastructure operator can:
\begin{itemize}
\item License the patent family before deploying PQ infrastructure --- cheapest moment to license, because no operational dependence on infringing infrastructure yet exists.
\item Avoid the patent family by designing the migration around it --- expensive, requires careful architectural choices the operator may not have visibility into yet.
\item Ignore and migrate without license --- creates litigation exposure once deployed.
\end{itemize}

The first option is overwhelmingly preferable for any rational operator. Lux's licensing pitch to a pre-migration target is therefore: ``You are required to migrate. Migration without license creates 14-year litigation exposure. License now at pre-migration pricing, structure your migration architecture with our IP cover, gain operational certainty.''

\subsection{The license-tier offering anchors institutional negotiations}

Tier-2 license (Lux-certified PQ rollup substrate) is the natural buying object for an institutional operator. The license bundles: (a) PQ patent license; (b) crypto-agility policy framework; (c) Lux Z-Chain substrate access; (d) Lux-certified compliance certification. The institutional buyer (BlackRock, Citi, JPMorgan) is buying not a patent license but a complete PQ-ready regulatory-compliant rollup substrate, with the patent license as one component of the bundle.

This converts the negotiation from defensive (``you're sued, settle'') to offensive (``you have a migration imperative, buy our platform''). The buyer's incentive is positive, and the deal size is platform-tier (5-10x patent-only deal size).

\subsection{The defensive coalition expands with each license}

Each Tier-1 ecosystem license, Tier-2 certified-rollup license, and Tier-3 external-operator license adds a participant to the Lux defensive coalition. The coalition obligates cross-license-back to Lux on any patents the licensee develops. Over the prosecution life of the family (2026-2046), the coalition accumulates licensee-back rights to potentially hundreds of PQ-blockchain patents filed by major operators. The Lux portfolio thus grows organically without further Lux prosecution investment.

\subsection{Anti-troll posture}

The PQ family's structure forecloses the standard anti-troll critique. Lux is not an entity sitting on a patent and asserting against operating companies; Lux is operating its own substrate (Z-Chain), licensing the substrate to ecosystem participants, and offering the patent as the technical-and-legal foundation of the substrate. The asymmetry between Lux (technology provider + IP holder) and pure-NPE (IP holder only) is structural and durable.

\section{Defense against acquirer / aggregator attacks}

The high-value structure of the family attracts predictable defensive responses from incumbents.

\subsection{Acquisition attempts}

Coinbase, Optimism Collective, Polygon, Anthropic, Google, IBM, Microsoft, and JPMorgan all have the resources and motivation to attempt acquisition of the patent family from Lux Industries. Acquisition removes the threat at one-time cost.

Lux defense: do not engage. The patent family is held within Lux Industries Inc. and is not transferable except through corporate-control change. Lux Industries' bylaws specifically restrict transfer of the patent family without supermajority shareholder approval. Acquisition offers are responded to with: ``Not for sale.''

\subsection{Defensive-aggregator attacks}

RPX, Allied Security Trust, Unified Patents, Open Invention Network, and the Crypto Open Patent Alliance (COPA) will attempt one or more of:

\begin{itemize}
\item Acquire the patent at distressed-asset pricing (typically 5-10\% of enforcement value).
\item File prophylactic IPRs to invalidate the patent before a complaint is served.
\item Recruit the patent into a defensive pool with cross-license obligations binding Lux to grant free licenses to pool members.
\item Lobby for regulatory or legislative changes that reduce enforceability.
\end{itemize}

Lux defense: bind the asset under exclusive litigation-finance arrangement that prevents asset sale. Block all communication channels from these entities. Move fast on filing first-wave complaints so that prophylactic IPR windows do not open in defendants' favor.

\subsection{Counter-invalidation alliance attacks}

Defendants may coordinate on a unified IPR strategy: file overlapping petitions, share prior-art research, retain shared expert witnesses. The unified attack increases per-defendant invalidity probability while reducing per-defendant cost.

Lux defense: vigorous claim construction at Markman in first cases to lock in patentee-favorable construction before unified IPR petitions are filed; engage former PTAB judges as consulting experts on IPR claim construction; pre-draft Markman briefs aligned with claim construction necessary to defeat unified anticipation/obviousness attacks.

\section{Conclusion}

The post-quantum patent family extends the strategic position created by US 11,803,853 from the present-day rollup infrastructure landscape to the future post-quantum infrastructure landscape, with no jurisdictional, primitive, or architectural escape route from the patent zone. Lux Industries' position as owner of both the original Nahmii patent and the PQ family provides a 13-year (Nahmii expiry 2039) plus 18-year (PQ family expiry 2044-2046, depending on prosecution timing) coverage window over the foundational architectural pattern of every operator-plus-validator-plus-execution-gate blockchain transaction-control system, classical or post-quantum.

The expected commercial value of the PQ family is harder to estimate than the original patent's because the migration window is in the future, but the relevant infrastructure-investment base is substantially larger. Conservative estimates place the post-2030 PQ-blockchain infrastructure-spending base at \$50B-\$200B annually across L1 consensus migration, L2 DA migration, cross-chain messaging migration, and institutional tokenization migration. A 0.25-1\% selective royalty applied to a meaningful fraction of this base supports a 10-20 year revenue program in the multi-billion-dollar range, with the substrate-licensing model providing additional platform-tier upside.

The strategic position is unique. It cannot be replicated by a competitor without filing comparable PQ-blockchain patents pre-priority date and prosecuting them through the same primitive-agnostic and architecture-agnostic framework. The Nahmii patent provided the foundation; the PQ family extends the foundation into the post-quantum era. Together they constitute the IP moat for Lux Industries' substrate-licensing business and the underlying enforcement asset for the multi-decade monetization program.

\section*{Cross-references}

\begin{itemize}
\item Nahmii patent acquisition paper: \texttt{../nahmii-patent-acquisition/nahmii-patent-acquisition.tex}
\item PQ master CIP application: \texttt{\textasciitilde/work/lux/patents/original-filings/2026-001-pq-transaction-control-master-cip/application.tex}
\item PQ provisionals 2026-002 through 2026-007: \texttt{\textasciitilde/work/lux/patents/original-filings/}
\item International filing strategy: \texttt{\textasciitilde/work/lux/patents/original-filings/international-filing-strategy.md}
\item Enforcement target ranking: \texttt{\textasciitilde/work/lux/patents/strategy/enforcement-targets.tex}
\item Corporate DAO stack analysis: \texttt{../corporate-dao-stack-analysis/corporate-dao-stack-analysis.tex}
\end{itemize}

\end{document}
